\section{Herní kazety a správa paměti (Mappers)}
\label{sec:mappers}

Na rozdíl od moderních konzolí, kde médium (CD, DVD, SD karta) slouží pouze jako pasivní úložiště
dat, byla herní kazeta pro NES (\textit{cartridge}) přímým rozšířením hardwaru konzole. Po zasunutí
do slotu, čipy na kazetě se mapovaly do adresního prostoru CPU i PPU. Existují stovky
zaregistrovaných mapperů, ale nejvyužívanější jsou NROM, MMC1, UNROM a MMC3.

\subsection{Struktura herní kazety a ROM obraz}

Každá standardní hra se skládá ze dvou základních typů paměti:

\begin{itemize}
    \item \textbf{PRG-ROM (Program ROM)} -- obsahuje kód hry a data pro procesor CPU. V adresním
        prostoru CPU je pro tuto paměť vyhrazen rozsah od \texttt{0x8000} do \texttt{0xFFFF} (celkem
        32~kB).
    \item \textbf{CHR-ROM (Character ROM)} -- obsahuje grafická data (dlaždice) pro procesor PPU. Tato
        data jsou mapována přímo do adresního prostoru grafického čipu. Pokud kazeta neobsahuje
        CHR-ROM, využívá tzv. \textbf{CHR-RAM}, do které si CPU nahraje grafiku z PRG-ROM při startu
        hry.
\end{itemize}

Digitální soubor, se kterým emulátor pracuje, se nazývá \textbf{ROM obraz} (nebo zkráceně „ROM“).
Ten vzniká procesem zvaným \textit{dumping}, kdy je obsah těchto fyzických čipů bit po bitu vyčten a
uložen do jednoho souboru. Aby však emulátor věděl, jak má s těmito daty naložit, musí být doplněna
metadata o hardwaru původní kazety.

\subsection{Formát iNES}

Standardem pro emulaci NES se stal formát .nes, vycházející ze specifikace iNES. Tento formát navrhl
Marat Fayzullin v roce 1996 pro svůj emulátor iNES a dodnes jde o nejrozšířenější způsob distribuce
NES her. Klíčovým prvkem je 16bajtová hlavička na začátku souboru (tabulka \ref{tab:ines_header}),
která popisuje konfiguraci hardwaru.

\begin{table}[htbp]
    \centering
    \begin{tabular}{llp{8cm}}
        \toprule
        \textbf{Byte} & \textbf{Význam} & \textbf{Popis} \\
        \midrule
        0--3 & Identifikátor & ASCII řetězec „NES“ následovaný hodnotou \texttt{0x1A}. \\
        4 & Velikost PRG-ROM & Počet 16~KiB bloků programové paměti. \\
        5 & Velikost CHR-ROM & Počet 8~KiB bloků grafické paměti. \\
        6 & Příznaky (Flags 6) & Obsahuje typ zrcadlení a spodní 4 bity čísla mapperu. \\
        7 & Příznaky (Flags 7) & Obsahuje horní 4 bity čísla mapperu a typ konzole. \\
        8--15 & Rezervováno & Většinou nevyužito (nebo použito pro novější standard NES 2.0). \\
        \bottomrule
    \end{tabular}
    \caption{Struktura hlavičky formátu iNES.}
    \label{tab:ines_header}
\end{table}

\subsection{Mechanismy bankování paměti (Mappers)}

Hlavním omezením procesoru R2A03 je jeho 16bitová adresní sběrnice, která mu umožňuje přímo
adresovat pouze 64 KiB paměti. Z tohoto prostoru je navíc pro hru vyhrazeno pouze 32 KiB, zbytek
zabírá RAM, registry periferií, zásobník a vektory přerušení. Jakmile hry začaly svou komplexností
tento limit přesahovat, byl na kazety přidáván hardware zvaný \textbf{mapper} (oficiálně \textit{Memory
Management Controllers} -- MMC).

Nejjednodušší mapper implementuje techniku přepínání bank (\textit{bank switching}). Fixní adresní
prostor CPU je dynamicky mapován na různé části mnohem větší fyzické paměti ROM na kazetě. Kromě
samotného mapování však pokročilé mappery nabízely i další funkce, jako např. vlastní zvukové
generátory navíc, PRG RAM (statická RAM paměť místo standardní ROM pamětí) nebo vlastní přerušení.
